\subsection{Messverfahren}
\xfigure{htbp}{width=0.8\textwidth}{Versuchsaufbau.png}{Versuchsaufbau\autocite{Skript_1}}{Versuchsaufbau\autocite{Skript_1}}
\paragraph{Geräte}
\begin{itemize}
    \item Kalorimeter mit Magnetrührer 
    \item Heizplatte mit Glasbecher 
    \item Stativ mit Drahthaken 
    \item 3 große Versuchskörper (Graphit, Aluminium, Blei) 
    \item 3 kleine Versuchskörper für den Versuch bei Stickstofftemperatur 
    \item Dewargefäß für flüssigen Stickstoff elektronisches Thermometer 
    \item elektronische Waage 
    \item Stoppuhr 
    \item Schutzbrille und Schutzhandschuhe 
\end{itemize}

\paragraph{Wasserwertbestimmung:}
Das leere Kalorimeter wird mit \(\qty{50}{\degreeCelsius}\) heißem Wasser befüllt und die Ausgleichstemperatur gemessen, die Masse des Wassers wird durch Wiegen mit und ohne Wasser bestimmt. Durch Aufzeichnen der Temperatur ergibt sich ein linearer Abfall durch Wärmeverlust, wird dieser extrapoliert, kann man die Ausgleichstemperatur bestimmen.
\paragraph{Messung von \textbf{\(c_{\mathrm{mol}}\)} bei rotem Temperaturbereich:}
In das wassergefüllte Kalorimeter werden auf \(\qty{100}{\degreeCelsius}\) erhitzte Probekörper eingetaucht und die Ausgleichstemperatur gemessen. Wasserdampfverluste werden vernachlässigt.

\paragraph{Messung von \textbf{\(c_{\mathrm{mol}}\):} im Temperaturbereich von flüssigem Stickstoff}
Durch Wiegen eines mit flüssigem Stickstoff befüllten Dewargefäßes vor und nach dem Eintauchen eines Probekörpers wird die Menge an verdampftem Stickstoff bestimmt.

\subsection{Messprotokoll}
\includegraphics[page=1,width=\textwidth]{V42.pdf}
\includegraphics[page=2,width=\textwidth]{V42.pdf}
\newpage
\includegraphics[page=3,width=\textwidth]{V42.pdf}
\newpage
\includegraphics[page=4,width=\textwidth]{V42.pdf}
\newpage
\includegraphics[page=5,width=\textwidth]{V42.pdf}